\documentclass[10pt]{article}
\usepackage[T1]{fontenc}
\usepackage[utf8]{inputenc}
\usepackage[french]{babel}
\usepackage{amsmath,amssymb,mathrsfs,mathpazo}
\usepackage{geometry}
\geometry{a4paper,margin=2cm}
\usepackage{graphicx,array,tabularx,tabularray,colortbl,multirow,multicol,booktabs,enumitem,pifont,ulem,fancybox,diagbox,cancel}
\usepackage{pstricks,pstricks-add,pst-plot,pst-tree,pst-node,pst-text,pst-func,pst-math,pst-bspline,pst-3dplot,pst-eucl,pst-solides3d}
\let\pscalculate\undefined
\usepackage{pst-bezier}
\usepackage{tikz,pgfplots,tkz-tab}
\usetikzlibrary{arrows,arrows.meta,patterns,calc,positioning,trees,decorations.markings,shapes.geometric,shapes.misc,shapes.arrows,intersections,decorations.pathmorphing,decorations.pathreplacing,backgrounds,shadows,babel}
\pgfplotsset{compat=1.18}
\usepgfplotslibrary{fillbetween}
\usepackage[np]{numprint}
\usepackage{listings,makecell,pas-tableur,twoopt}
\let\barre\undefined
\usepackage[pstricks]{tablvar}
\definecolor{darkgreen}{RGB}{0,100,0}
\providecommand{\red}{\color{red}}
\providecommand{\blue}{\color{blue}}
\providecommand{\green}{\color{green}}
\providecommand{\black}{\color{black}}
\providecommand{\white}{\color{white}}
\providecommand{\gray}{\color{gray}}
\providecommand{\cyan}{\color{cyan}}
\providecommand{\magenta}{\color{magenta}}
\providecommand{\yellow}{\color{yellow}}
\providecommand{\violet}{\color{violet}}
\providecommand{\orange}{\color{orange}}
\providecommand{\euro}{\texteuro}
\providecommand{\EUR}{\texteuro}
\providecommand{\ieme}{\textsuperscript{e}}
\providecommand{\degres}{\ensuremath{^\circ}}
\providecommand{\degree}{\ensuremath{^\circ}}
\providecommand{\vv}[1]{\overrightarrow{#1}}
\providecommand{\Coord}[1]{#1}
\providecommand{\decosix}{}
\providecommand{\starredbullet}{\ensuremath{\bullet}}
\pagestyle{empty}
\setlength{\parindent}{0pt}
\usepackage{adjustbox}\DeclareUnicodeCharacter{2236}{:}\DeclareUnicodeCharacter{2190}{\ensuremath{\leftarrow}}
\providecommand{\R}{\mathbb{R}}
\providecommand{\N}{\mathbb{N}}
\providecommand{\Z}{\mathbb{Z}}
\providecommand{\Q}{\mathbb{Q}}
\providecommand{\D}{\mathbb{D}}
\providecommand{\C}{\mathbb{C}}
\providecommand{\e}{{\mathrm{e}}}
\providecommand{\vect}[1]{\overrightarrow{#1}}
\providecommand{\vectt}[1]{\overrightarrow{\mathrm{#1}}}
\providecommand{\vvt}[1]{\overrightarrow{\mathrm{#1}}}
\providecommand{\np}[2][]{{#2~#1}}
\providecommand{\vv}[1]{{\overrightarrow{#1}}}
\renewcommand{\d}{{\mathrm{d}}}
\providecommand{\degres}{{^\circ}}
\providecommand{\degre}{{^\circ}}
\providecommand{\blue}{\color{blue}}
\providecommand{\red}{\color{red}}
\providecommand{\enskip}{\quad}
\providecommand{\llbracket}{\lbrack\!\lbrack}
\providecommand{\rrbracket}{\rbrack\!\rbrack}
\providecommand{\euro}{{\text{€}}}
\providecommand{\no}{{\text{n°}}}
\providecommand{\interval}[2]{{[#1;#2]}}
\providecommand{\intervFO}[2]{{[#1;#2[}}
\providecommand{\dotfill}{\ldots}

\providecommand{\ds}{\displaystyle}
\providecommand{\textding}[1]{#1}
\providecommand{\Oij}{\(\mathrm{O};\vec{\imath},\vec{\jmath}\)}
\providecommand{\Oijk}{\(\mathrm{O};\vec{\imath},\vec{\jmath},\vec{k}\)}
\providecommand{\Ouv}{\(\mathrm{O};\vec{u},\vec{v}\)}
\providecommand{\intervFF}[2]{\ensuremath{[#1;#2]}}
\renewcommand{\euro}{\ifmmode\text{\texteuro}\else\texteuro\fi}
\DeclareUnicodeCharacter{00B0}{\ifmmode^{\circ}\else\textdegree\fi}
\DeclareUnicodeCharacter{20AC}{\ifmmode\text{\texteuro}\else\texteuro\fi}
% Énoncés : Éducation nationale. Transcription et corrigés : APMEP, Denis Vergès et contributeurs. Découpage et classement : CoopMaths, Sébastien Lozano, Éric Elter et contributeurs. Adaptation de la présentation : Maths Explorer.
% CC BY-SA : https://coopmaths.fr/www/about/
% https://coopmaths.fr/alea/static/e3c/2021/tex/e3c_2021_01_specimen2_2_cor.tex
\begin{document}
\noindent\textit{APMEP / CoopMaths. CC BY-SA. Présentation adaptée par Maths Explorer.}\par\medskip
\bigskip

Une collectivité locale octroie une subvention de \np{116610}~\euro{} pour le forage d'une nappe d'eau souterraine. Une entreprise estime que le forage du premier mètre coûte 130~\euro; le
forage du deuxième mètre coûte 52~\euro{} de plus que celui du premier mètre ; le forage du
troisième mètre coûte 52~\euro{} de plus que celui du deuxième mètre, etc.

Plus généralement, le forage de chaque mètre supplémentaire coûte 52~\euro{} de plus que celui
du mètre précédent.

Pour tout entier $n$ supérieur ou égal à 1, on note: $u_n$ le coût du forage du $n$-ième mètre en
euros et $S_n$ le coût du forage de $n$ mètres en euros; ainsi $u_1=130$.

\medskip

\begin{enumerate}
\item  $u_2=130+52=182$ et $u_3=182+52= 234$.
\item La suite $\left(u_n\right)$ est arithmétique de premier terme $u_1=130$ et de raison $r=52$.



On en déduit que pour tout $n\geqslant 1$, on a: $u_n=u_1+(n-1)\times r$ donc $u_n=130+52(n-1)$.

\item $S_2= u_1+u_2 = 130+182=312$ et $S_3=u_1+u_2+u_3 = S_2+u_3=312+234=546$.

\item Afin de déterminer le nombre maximal de mètres que l'entreprise peut forer avec la subvention qui est octroyée, on considère la fonction Python suivante:

\begin{center}
\fbox{\tt
\begin{tabular}{l}
def nombre\_metre(S) :\\
\quad C = 130\\
\quad n = 1\\
\quad while C < S :\\
\qquad C = C + ...\\
\qquad n = n + 1\\
 \quad return n
\end{tabular}
}
\end{center}

On veut compléter cet algorithme de sorte que l'exécution de la fonction \verb!nombre_metre(S)! renvoie le nombre maximal de mètres que l'entreprise peut forer avec la subvention octroyée. 

La variable \texttt{C} contient le coût du forage de \texttt{n} mètres.

Pour tout $n\geqslant 1$, on a $S_{n+1}= u_1+u_2+\cdots + u_n + u_{n+1} = S_n+u_{n+1}$; or $u_n=130+52\left (n-1\right )$ donc $u_{n+1} =130+52n$. On a donc $S_{n+1} = S_n + 130 + 52n$.

Il faut donc ajouter $130+52n$ à la variable  \texttt{C} à chaque tour de boucle.

\begin{center}
\fbox{\tt
\begin{tabular}{l}
def nombre\_metre(S) :\\
\quad C = 130\\
\quad n = 1\\
\quad while C < S :\\
\qquad C = C + \blue 130 + 52*n\\
\qquad n = n + 1\\
 \quad return n
\end{tabular}
}
\end{center}

\item On admet que, pour tout entier naturel non nul, $S_n=26n^2 +104n$. 

On cherche le plus grand entier $n$ tel que $S_n \leqslant \np{116610}$.

On résout l'équation $26n^2 + 104n - \np{116610}=0$.

$\Delta = 104^2 - 4 \times 26 \times (-\np{116610}) = \np{12138256} = \np{3484}^2$ 

Deux solutions: $n'=\dfrac{-104+\np{3484}}{2\times 26}=65$ et $n''=\dfrac{-104-\np{3484}}{2\times 26}=-69<0$.

Avec $\np{116610}$~\euro{}, on pourra creuser 65 mètres.

\end{enumerate}

\vspace{0.5cm}

\end{document}
